\section{Sesgos, ética y privacidad}

\subsection{Sesgo de muestreo}

La fuente contiene exactamente 1.000 filas por cada uno de 114 géneros. Este balance facilita
comparaciones, pero no representa la prevalencia natural de canciones ni de reproducciones en
Spotify. Además, las canciones multigénero aparecen repetidas en el origen, por lo que estadísticas
agregadas ponderan más a pistas con varias etiquetas. La consolidación corrige el doble conteo por
canción, pero no reconstruye un muestreo aleatorio del catálogo completo.

\subsection{Sesgo temporal}

\code{popularity} depende de reproducciones y recencia \citep{spotify_track_api}. Al faltar fecha de
extracción y lanzamiento, no es posible distinguir canciones nuevas en crecimiento de clásicos
estables ni evaluar evolución. Un modelo entrenado en esta fotografía puede sufrir \emph{concept
drift}: la relación entre atributos y popularidad cambia con tendencias, reglas de plataforma y
audiencias.

\subsection{Variables omitidas}

No se observan inversión de marketing, exposición en playlists, viralidad, popularidad previa del
artista, colaboraciones, región, fecha de lanzamiento, estacionalidad ni exposición externa. Estas
variables podrían explicar simultáneamente atributos y popularidad, de modo que las correlaciones de
audio no deben interpretarse como efectos. Su ausencia también fija un techo razonable a la capacidad
predictiva.

\subsection{Sesgo y memorización por artista}

Hay 31.437 valores de \code{artists}. Una partición aleatoria puede dejar canciones del mismo artista
en entrenamiento y prueba; el modelo reconocerá reputación histórica y aparentará generalizar mejor
de lo que lo haría con artistas nuevos. Para estimar ese escenario se necesita validación agrupada por
artista o una evaluación separada de artistas no vistos.

\subsection{Sesgo por género y multietiqueta}

Una canción puede contribuir a varias etiquetas en análisis expandidos. Comparar géneros como grupos
excluyentes sería incorrecto, y un multi-hot puede aprender asociaciones del diseño de muestreo. Las
métricas por género deben informar soporte, permitir pertenencia múltiple y evitar concluir que una
etiqueta determina el desempeño.

\subsection{Riesgos éticos}

Una predicción puede reforzar artistas ya populares, favorecer estilos dominantes, reducir
oportunidades de artistas nuevos y homogeneizar decisiones creativas. También puede confundirse
probabilidad estimada con mérito artístico o usarse como regla automática de inversión. Una
correlación como la de \code{instrumentalness} no autoriza recomendar eliminar elementos
instrumentales: refleja esta muestra y puede estar mediada por géneros y exposición.

Principios de uso responsable:

\begin{itemize}
\item usar el modelo como apoyo a decisión humana y documentar su propósito;
\item comunicar incertidumbre, datos de entrenamiento y límites de generalización;
\item revisar errores y cobertura por género, condición explícita y artistas vistos/no vistos;
\item monitorear deriva y recalibrar solo con datos trazables;
\item evitar asignaciones automáticas de presupuesto o promoción de alto impacto;
\item preservar diversidad creativa y habilitar revisión de decisiones.
\end{itemize}

\subsection{Privacidad}

El dataset contiene metadatos de canciones, nombres públicos de artistas y atributos musicales. No
incluye directamente identidad, ubicación, perfiles ni historial de escucha de oyentes; el riesgo de
privacidad individual es relativamente bajo frente a datos conductuales. Aun así, se mantiene
trazabilidad, minimización y respeto por las condiciones de uso de la fuente.

Si una etapa futura integra usuarios, playlists personales, ubicaciones o eventos de escucha, cambia
materialmente el riesgo. Serían necesarios propósito y base legal definidos, minimización, controles
de acceso, retención limitada, agregación o seudonimización, revisión de reidentificación y una
evaluación de impacto antes del uso.
